% =====================================================================
%  MATH3030 Project Report
%  Build:  latexmk          (uses .latexmkrc -> XeLaTeX + biber)
% =====================================================================
\documentclass[a4paper, 11pt]{article}

% --- 1. Packages that do things (maths, figures, refs) ---------------
% =====================================================================
%  preamble.tex - functional packages only.
%  Nothing here should change how the document *looks*; that belongs in
%  themes/. Keep this file the same when you switch themes.
% =====================================================================

% --- Maths ------------------------------------------------------------
\usepackage{amsmath, amssymb, amsthm}
\usepackage{mathtools}

% --- Figures & tables -------------------------------------------------
\usepackage{graphicx}
\graphicspath{{figures/}}
\usepackage{subcaption}   % side-by-side figures: \begin{subfigure}
\usepackage{booktabs}     % nice tables: \toprule \midrule \bottomrule
\usepackage{float}
\usepackage{comment}      % \begin{comment} ... \end{comment} to hide blocks

% --- Bibliography -----------------------------------------------------
\usepackage[style=authoryear, backend=biber, maxcitenames=2]{biblatex}
%  Use \textcite{turing1952} -> "Turing (1952)"
%      \parencite{turk1991}  -> "(Turk 1991)"

% --- Links & cross-references (load hyperref late, cleveref last) -----
\usepackage{hyperref}
\usepackage[capitalise, noabbrev]{cleveref}
%  Use \cref{eq:gray-scott} -> "Equation 1", \cref{fig:spots} -> "Figure 2"

% --- Equation numbering by section: (2.1), (2.2), ... ------------------
\numberwithin{equation}{section}

% --- Theorem-like environments ----------------------------------------
\theoremstyle{definition}
\newtheorem{definition}{Definition}[section]
\newtheorem{example}[definition]{Example}
\theoremstyle{plain}
\newtheorem{theorem}[definition]{Theorem}
\newtheorem{proposition}[definition]{Proposition}
\theoremstyle{remark}
\newtheorem*{remark}{Remark}

% --- Drafting helper: \todo{...} shows in red; disable before submitting
\usepackage{xcolor}
\newcommand{\todo}[1]{\textcolor{red}{[TODO: #1]}}
% \renewcommand{\todo}[1]{}   % <- uncomment to hide all TODOs


% --- 2. Notation shortcuts -------------------------------------------
% =====================================================================
%  macros.tex - notation shortcuts. Define a symbol once here so it is
%  consistent everywhere, and can be changed in one place.
% =====================================================================

\newcommand{\R}{\mathbb{R}}
\newcommand{\diff}{\mathop{}\!\mathrm{d}}          % upright d: \diff t
\newcommand{\pd}[2]{\frac{\partial #1}{\partial #2}} % \pd{u}{t}
\newcommand{\lap}{\nabla^2}                          % Laplacian
\newcommand{\Jac}{J}                                 % Jacobian matrix
\DeclareMathOperator{\tr}{tr}                        % trace
% \det already exists in LaTeX

% Gray-Scott parameters (change symbols here if you switch convention)
\newcommand{\GSF}{F}
\newcommand{\GSk}{k}
\newcommand{\Du}{D_u}
\newcommand{\Dv}{D_v}


% --- 3. THEME: everything about how the document LOOKS ----------------
%  To change the look, point this at a different file in themes/.
%    themes/basic.tex  - plain LaTeX defaults (current)
%    themes/serif.tex  - Georgia + tighter margins, like your project plan
% =====================================================================
%  themes/basic.tex - a plain, safe starting look.
%  A theme file controls ONLY appearance: page size/margins, fonts,
%  heading style, spacing, colours, header/footer.
%  To make a new theme: copy this file, edit it, and change the
%  \input{themes/...} line in main.tex.
% =====================================================================

% --- Page layout ------------------------------------------------------
\usepackage[margin=2.5cm]{geometry}

% --- Fonts ------------------------------------------------------------
%  Default Latin Modern (works with pdfLaTeX and XeLaTeX).

% --- Paragraphs: space between, no indent ------------------------------
\usepackage[skip=6pt plus 1pt]{parskip}

% --- Links: no coloured boxes ------------------------------------------
\hypersetup{
  colorlinks = true,
  linkcolor  = black,
  citecolor  = black,
  urlcolor   = blue!50!black,
}


% --- 4. Bibliography --------------------------------------------------
\addbibresource{references.bib}

% --- Title information -------------------------------------------------
\title{Nonlinear Dynamics in Computer Graphics:\\
       Reaction--Diffusion Systems in Procedural Generation}
\author{Ragav Srinivasan}
\date{\today}

\begin{document}

\maketitle

\begin{abstract}
	We start by ti
\end{abstract}

\section{Introduction}
\label{sec:introduction}
% Structure 
\subsection{Structure of this report}
% Describe each section in their own line:
\Cref{sec:background} gives the background on reaction--diffusion systems and the phenomena they
model. We then discuss a famous result of Alan Turing on reaction--diffusion systems, and finally
focus on a simple reaction--diffusion system called the Gray--Scott model.
\Cref{sec:simulations} is my original contribution: a stability analysis of the Gray--Scott model
and a numerical simulation of it.
\Cref{sec:conclusion} gives the conclusion and possible extensions of this project.
\section{Background}
\label{sec:background}



\subsection{Procedural Generation}
\label{sec:procedural}

\begin{itemize}
	\item Procedural Generation is a technique in Computer Graphics that generates smooth textures and landscapes in real time. Artists and Developers can create \emph{lifelike} 3D terrains with mountains, rivers and vegetation without the need to design or program them manually. To the player, these terrains will generate on the fly no pre-processing required.
	\item The most famous use case of Procedural Generation can be found in Minecraft which makes heavy use of procedural techniques to generate worlds and cave systems.
	\item Procedural generation is a broad technique that encompasses many algorithms. However, we will particularly looking at one specific algorithm: \emph{Perlin Noise}.
\end{itemize}


% Origin of Perlin Noise 
Perlin Noise was devised by Ken Perlin at The Mathematical Animation Group while working for the 1982 Disney Sci-fi film \emph{Tron}. Perlin's noise function still sees widespread use even in modern day games.

% Mathematics behind Perlin noise
\todo{Maths behind perlin noise}
\begin{align*}
	noise(x) = x
\end{align*}
% Use cases of Perlin noise 

\begin{itemize}
	\item Perlin noise is a continuous function that uses
\end{itemize}

%  Limitations of noise functions:
Textures created by Perlin Noise and other noise functions are random and lack structure \cite{lagae2010}. This is a useful feature for landscapes but if we wanted a zebra whose stripes have self organized patterns, then a different approach is in order.


\subsection{Reaction Diffusion Systems}
\label{sec:rd-systems}

% Segway to Reaction Diffusion Systems 

This brings us to the heart of this report, \emph{Reaction Diffusion Systems}. These are Nonlinear dynamical systems originating from the world of Chemistry that models the intereactions between multiple chemical susbtances and their diffusivity with resepct to space and time.
\\
\emph{Diffusion} refers to the spread of chemical substances between low and high concentration. It turns out adding an element of diffusion leads to some very interesting pattern formation.
\\
The General Form of Reaction Diffusion Systems is given by:

\begin{equation} \label{eq:rd-general}
	\pd{\mathbf{q}}{t} = D \nabla^2 \mathbf{q} + R(\mathbf{q}),
\end{equation}


% Explanation of the system
\begin{itemize}
	\item The equation is \emph{dimensionless}
	\item $\mathbf{q}$ is a vector of unknown vector functions representing the concentrations of an array of chemical substances.
	\item $D$ gives the diffusivity for each chemical
	\item $R(\mathbf{q}$) is another vector functions that describes all the chemical reactions.
\end{itemize}

There are several systems that fit the criteria of Reaction Diffusion Systems. Notably, the activator eqution, some russian equation, etc.


\subsection{Turing patterns}
\label{sec:turing-patterns}
\todo{Define Turing patterns} \\
Turing patterns are self organized patterns that resembles structures found in nature. \\
\todo{Cite turing's 1952 discovery on chemical morphogenesis}


\subsection{The Gray-Scott model}
\label{sec:gray-scott}

The Gray Scott model is a simple reaction diffusion system that models the interaction between chemicals $U$ and $V$ based on the following rule:

\begin{align*}
	U + 2V \to 3V \\
	V \to P
\end{align*}


% General form of the Gray Scott Model
\begin{align*}
	\frac{\partial U}{\partial t} & = -UV^2 + D_u \nabla^2 + f(1 - U) \\
	\frac{\partial V}{\partial t} & = UV^2 + D_v \nabla^2 + (k + f)V
\end{align*}
% Explaination of the terms 
\begin{itemize}
	\item $D_u$ and $D_v$ is the diffusivity of chemicals $U$ and $V$.
	\item $U$ and $V$ are the unknown vector functions representing the concentrations of the respective chemicals in space and time.
	\item $f$ is the feed rate of chemical $U$ for $V$
	\item $k$ is the decay rate for chemical $V$.
\end{itemize}

% What we are going to do with this model 
The broader goal of this project is to explain pattern formation rising from Reaction Diffusion Systems but the narrower goal is to find which parameter values of $F$ and $k$. Therefore, In the end we will build a parameter map of the Gray Scott Model that can identify which patterns we get.


\section{Original Contribution}
\label{sec:simulations}

\subsection{Stability analysis}
\label{sec:stability}


% Definition of Stability Analysis
A stability analysis gives us a guess on the behaviour of the system. For instance, a system is said to be stable if solutions move away from the equillibrium point, otherwise an unstable system is when solutions move away from the equillibrium point.


Nonlinear Dynamical Systems are hard to solve analytically.




% Why are we studying the stability analysis for the gray scott model
The point of this stability analysis is to show how adding a diffusion effect can destabilize a system. This is also called as Turing instability.

% How do we do a Stability Analysis:
In this case, we will be doing a stability analysis twice. The first part will prove the Gray Scott Model is stable when we drop diffusion. In the second part we will show that the Gray Scott Model's turing unstable after adding the diffusion constant.


%% Maths for Stability analysis 
1. Setup the Gray Scott system  as follows:
\begin{align*}
	\dot{U} = -UV^2 + D_u \nabla^2 + F(1 - U) \\
	\dot{V} = UV^2 + D_v \nabla^2 + (F + k)V
\end{align*}

2. What are the equillibrium points:
From inspection, there are two $(0, 0)$ and $(1, 0)$.

3. Jacobian of the gray-scott model:
\begin{align*}
	J = \begin{pmatrix}
		    \partial_{u} \dot{U} & \partial_v  \dot{U}  \\
		    \partial_{u} \dot{V} & \partial_{v} \dot{V}
	    \end{pmatrix}
	=
	\begin{pmatrix}
		- V^2 - F & -2UV     \\
		V^2       & (F + k)V
	\end{pmatrix}
\end{align*}

4. Substitute equillibrium points into jacobian

First equillibrium point:
\begin{align*}
	J(0, 0) = \begin{pmatrix}
		          -F & 0       \\
		          0  & (F + k)
	          \end{pmatrix}
\end{align*} \\

$p = Tr(J(0,0)) = k$  and  $q = \det(J(0,0)) = -F(F + k)$ \\
what is $p^2 - q$? \\
\todo{Write your stability analysis} Therefore, the system is \emph{stable}.



Substituting the second equillibrium point yields:
\begin{align*}
	J(1, 0) = \begin{pmatrix}
		          -F & 0 \\
		          0  & 0
	          \end{pmatrix}
\end{align*}
$Tr(J(1,0)) = -F$  and  $\det(J(1,0)) = 0$ \\
Therefore, the system is \emph{stable}.



Part 2: We do the same analysis with the diffusion terms added back?

1. Equillibrium points
2. Jacobian
3. Tr J and Det J of the jacobian
4. Is it stable or unstable.

\subsection{Simulation of the Gray Scott Model}
\label{sec:numerical}
% Movtiation for doing a numerical simulation
\subsection{Texture Mapping in Blender}
\todo{I will do this section if I have time, for now I will just do the parameter map and stability analysis.}



\todo{Method details, figures and discussion.}
\section{Conclusion}
\label{sec:conclusion}
\todo{Summarise what the stability analysis predicts, how well the simulation matches it, and what
	this means for someone choosing $F$ and $k$ to generate textures. Possible extensions: varying the
	diffusion ratio $D_u/D_v$, comparing the predicted and simulated spot spacing, and running the
	simulation on a 3D surface as in \textcite{turk1991}.}

\printbibliography


\section{Appendix}
\end{document}

